\documentclass[10pt,a4paper]{article}
\usepackage{amsmath,amssymb,amsthm,graphicx}
\usepackage[shortlabels]{enumitem}
\usepackage{amscd,mathtools,mathrsfs,graphicx,xcolor}
\usepackage[all,cmtip]{xy}
\let\colonsim\relax
\let\colonapprox\relax
\usepackage{colonequals}
\usepackage[a4paper,margin=25mm]{geometry}
\usepackage[colorlinks=true,linkcolor=blue,citecolor=blue,urlcolor=blue]{hyperref}
\mathtoolsset{showonlyrefs=true}
\numberwithin{equation}{section}
\emergencystretch=3em
\newcommand{\doi}[1]{\href{https://doi.org/#1}{doi:\nolinkurl{#1}}}
\newcommand{\arXiv}[1]{\href{https://arxiv.org/abs/#1}{arXiv:\nolinkurl{#1}}}
\newtheorem{theorem}{Theorem}[section]
\newtheorem{lemma}[theorem]{Lemma}
\newtheorem{conjecture}[theorem]{Conjecture}
\newtheorem{corollary}[theorem]{Corollary}
\newtheorem{proposition}[theorem]{Proposition}

\newtheorem{thmAlph}{Theorem}
\renewcommand*{\thethmAlph}{\Alph{thmAlph}}
\newtheorem{corAlph}{Corollary}[thmAlph]

\newtheorem{thmIntro}{Theorem}
\newtheorem{conjIntro}{Conjecture}

\theoremstyle{definition}
\newtheorem{definition}[theorem]{Definition}
\newtheorem{condition}[theorem]{Condition}
\newtheorem{assumption}[theorem]{Assumption}

\theoremstyle{remark}
\newtheorem{rmk}[theorem]{Remark}


%%%%%%%%%%%%%%%%%%%%%%%%%%%%%%%%%%%%%
% Define any new command you need
%%%%%%%%%%%%%%%%%%%%%%%%%%%%%%%%%%%%%


\makeatletter
\newcommand{\sigmaop}[1]{\mathop{\mathpalette\@sigmaop{#1}}\slimits@}
\newcommand{\@sigmaop}[2]{%
  \vphantom{\sum}%
  \sbox\z@{$\m@th#1\sum$}%
  \dimen@=\ht\z@ \advance\dimen@\dp\z@
  \dimen\tw@=\wd\z@
  \ifx#1\displaystyle\dimen@=.9\dimen@\fi
  \ooalign{%
    \hidewidth
    $\vcenter{\hbox{$\m@th#1#2$\kern.3\dimen\tw@}%
     \ifx#1\scriptstyle\kern-.25ex\fi}$\hidewidth\cr
    $\vcenter{\hbox{%
      \resizebox{!}{\dimen@}{$\m@th\boxtimes$}%
    }\ifx#1\scriptstyle\kern-.25ex\fi}$\cr
  }%
}
\makeatother


%================Lists

\newlist{equivlist}{enumerate}{2}
\setlist[equivlist,1]{label={\rm(\roman*)}, ref={\rm\roman*}}

\newlist{assertionlist}{enumerate}{2}
\setlist[assertionlist,1]{label={\rm(\arabic*)}, ref={\rm\arabic*}}

\newlist{Alist}{enumerate}{2}
\setlist[Alist,1]{label={\rm(\Alph*)}, ref={\rm\Alph*}}

\newlist{alist}{enumerate}{2}
\setlist[alist,1]{label={\rm(\alph*)}, ref={\rm\alph*}}

%===================

\newcommand{\GZip}{\mathop{\text{$G$-{\tt Zip}}}\nolimits}

\newcommand{\pha}{\mathsf{pHa}}
\newcommand{\zipsf}{\mathsf{zip}}
\newcommand{\GSsf}{\mathsf{GS}}
\newcommand{\lwsf}{\mathsf{lw}}
\newcommand{\hwsf}{\mathsf{hw}}
\newcommand{\low}{\mathsf{low}}
\newcommand{\high}{\mathsf{high}}

\newcommand{\flag}{\mathsf{flag}}
\newcommand{\orb}{\mathsf{orb}}
\newcommand{\unip}{\mathsf{unip}}
\newcommand{\Lmin}{\text{$L$}\mathsf{-min}}

\newcommand{\GFqrZip}{\mathop{\text{$G_{\FF_{q^r}}$-{\tt Zip}}}\nolimits}
\newcommand{\GjZip}{\mathop{\text{$G_j$-{\tt Zip}}}\nolimits}
\newcommand{\GtilZip}{\mathop{\text{$\widetilde{G}$-{\tt Zip}}}\nolimits}

\newcommand{\GaZip}{\mathop{\text{$G_\alpha$-{\tt Zip}}}\nolimits}
\newcommand{\SLZip}{\mathop{\text{$\SL_2$-{\tt Zip}}}\nolimits}
\newcommand{\HZip}{\mathop{\text{$H$-{\tt Zip}}}\nolimits}
\newcommand{\LzZip}{\mathop{\text{$L_0$-{\tt Zip}}}\nolimits}
\newcommand{\GoneZip}{\mathop{\text{$G_1$-{\tt Zip}}}\nolimits}
\newcommand{\GtwoZip}{\mathop{\text{$G_2$-{\tt Zip}}}\nolimits}
\newcommand{\GspZip}{\mathop{\text{$\GSp(2g)$-{\tt Zip}}}\nolimits}
\newcommand{\GpZip}{\mathop{\text{$G'$-{\tt Zip}}}\nolimits}
\newcommand{\GF}{\mathop{\text{$G$-{\tt ZipFlag}}}\nolimits}
\newcommand{\HF}{\mathop{\text{$H$-{\tt ZipFlag}}}\nolimits}

\newcommand{\VB}{\mathfrak{VB}}

%================================================================

\newcommand{\Acal}{{\mathcal A}}
\newcommand{\Bcal}{{\mathcal B}}
\newcommand{\Ccal}{{\mathcal C}}
\newcommand{\Fcal}{{\mathcal F}}
\newcommand{\Lcal}{{\mathcal L}}
\newcommand{\Ocal}{{\mathcal O}}
\newcommand{\Ucal}{{\mathcal U}}
\newcommand{\Vcal}{{\mathcal V}}
\newcommand{\Xcal}{{\mathcal X}}
\newcommand{\Zcal}{{\mathcal Z}}
\newcommand{\pfr}{{\mathfrak p}}
\newcommand{\Afr}{{\mathfrak A}}
\newcommand{\Dfr}{{\mathfrak D}}
\newcommand{\Ifr}{{\mathfrak I}}

\renewcommand{\AA}{\mathbb{A}}
\newcommand{\CC}{\mathbb{C}}
\newcommand{\FF}{\mathbb{F}}
\newcommand{\GG}{\mathbb{G}}
\newcommand{\NN}{\mathbb{N}}
\newcommand{\QQ}{\mathbb{Q}}
\newcommand{\RR}{\mathbb{R}}
\newcommand{\XX}{\mathbb{X}}
\newcommand{\ZZ}{\mathbb{Z}}

\newcommand{\Bscr}{{\mathscr B}}
\newcommand{\Gscr}{{\mathscr G}}
\newcommand{\Lscr}{{\mathscr L}}
\newcommand{\Pscr}{{\mathscr P}}
\newcommand{\Sscr}{{\mathscr S}}
\newcommand{\Tscr}{{\mathscr T}}

%%%%%%%%%%%%%%%%%%%%%%%%%%%%%%%%%%%%%%%%

\newcommand{\cent}{{\rm Cent}}
\newcommand{\gm}{\mathbf G_{\textnormal{m}}}
\newcommand{\gmq}{\mathbf G_{\textnormal{m},\QQ}}
\newcommand{\fp}{\mathbf F_p}

\DeclareMathOperator{\codim}{codim}

\renewcommand{\div}{\operatorname{div}}

%DeclareMathOperator-----------------------------------------
\DeclareMathOperator{\ad}{ad}
\DeclareMathOperator{\card}{Card}

\newcommand{\dR}{{\rm dR}}

\DeclareMathOperator{\Gal}{Gal}
\DeclareMathOperator{\gal}{Gal}
\DeclareMathOperator{\Hom}{Hom}
\DeclareMathOperator{\rank}{rank}
\DeclareMathOperator{\Span}{Span}
\DeclareMathOperator{\Stab}{Stab}
\DeclareMathOperator{\pr}{pr}
\DeclareMathOperator{\Ker}{Ker}
\DeclareMathOperator{\Rep}{Rep}
\DeclareMathOperator{\res}{Res}
\DeclareMathOperator{\Sbt}{Sbt}
\DeclareMathOperator{\Sh}{Sh}
\DeclareMathOperator{\spec}{Spec}

\DeclareMathOperator{\zip}{zip} % not an operator? 

\DeclareMathOperator{\GS}{GS}
\DeclareMathOperator{\SL}{SL}
\DeclareMathOperator{\GL}{GL}
\DeclareMathOperator{\GSp}{GSp}
\DeclareMathOperator{\Sp}{Sp}
\DeclareMathOperator{\U}{U}
\DeclareMathOperator{\GU}{GU}



%--------------------------------------------------------
\newcommand{\shgx}{\Sh(\mathbf G, \mathbf X)}
\newcommand{\egx}{E(\mathbf G, \mathbf X)}
\newcommand{\gx}{(\mathbf G, \mathbf X)}

\newcommand{\gofaf}{\mathbf G(\mathbf A_f)}

\newcommand{\id}{{\rm Id}}

%Theorem etc. abbreviations--------------------------------------------------------------
\newcommand{\loccit}{{\em loc.\ cit. }} % check spacing 
\newcommand{\loccitn}{{\em loc.\ cit.}}
\newcommand{\opcitn}{{\em op.\ cit.}}
\newcommand{\cf}{{\em cf. }} 

%-----------------------------------------------------

\DeclareMathOperator{\diag}{diag}
\DeclareMathOperator{\fil}{Fil}

\newcommand{\reg}{{\rm reg}}


\DeclareMathOperator{\GSpin}{GSpin}
\DeclareMathOperator{\Spin}{Spin}
\DeclareMathOperator{\SO}{SO}
\DeclareMathOperator{\Fil}{Fil}

\DeclareMathOperator{\Ha}{Ha}
\DeclareMathOperator{\type}{type}
\DeclareMathOperator{\Norm}{Norm}
\DeclareMathOperator{\Ind}{Ind}
\DeclareMathOperator{\cara}{char}

\DeclareMathOperator{\Res}{Res}
\DeclareMathOperator{\Flag}{Flag}
\DeclareMathOperator{\Aut}{Aut}
\DeclareMathOperator{\ev}{ev}
\DeclareMathOperator{\Weil}{Weil}

\newcommand{\relmiddle}[1]{\mathrel{}\middle#1\mathrel{}}

%% 
\newcommand{\lra}{\longrightarrow}
\def\lowsim{\vbox to 0pt{\vss\hbox{$\scriptstyle\sim$}\vskip-1.6pt}}

\def\wmu{\raisebox{-.3em}{\scriptsize{$\widetilde{\mu}$}}}

\newcommand{\supth}[1]{\ensuremath{#1^{\mathrm{th}}}}
\usepackage{cleveref}
\crefname{originalproposition}{original proposition}{original propositions}
\Crefname{originalproposition}{Original proposition}{Original propositions}
\makeatletter
\newlabel{prop-charL}{{3.8}{source}{}{sourcepointer-1}{}}
\newlabel{prop-charL@cref}{{[originalproposition][0][]3.8}{[1][0][]source}{}{}{}}
\makeatother
\title{Highest-weight norm extension criterion for general G-zips (Theorem 4 / Proposition 3.9)}
\author{Naoki Imai, Jean-Stefan Koskivirta}
\date{arxiv:2211.16207; CC BY-SA 4.0}
\begin{document}
\maketitle
\noindent\textit{Editorial scope.} One entire norm-extension iff criterion and complete current proof; no mathematical repair or formal admission.\par\medskip
\noindent\textit{Allowed standard prerequisites (editorial).}\begin{itemize}
\item Zip orbit parametrization and finite stabilizer. [PWZ11] Theorems 7.5, 11.2, 11.3, 11.5; [KW18] Lemma 3.2.1; [IK21] Corollary 3.2.3.. Original Section 2 explicitly states orbit and invariant descriptions; these are earlier external results.
\item Root-curve boundary test and Frobenius factorization. [IK21] Automorphic vector bundles on the stack of G-zips, Forum Math. Sigma 9 (2021), e37, DOI 10.1017/fms.2021.32, Section 3.1 and Proposition 3.1.4; [Kos19] Section 3.2.. Earlier geometric extension test and root-curve factorization are explicitly attributed at source lines 1043, 1059, 1070 and 1082. The current general-parabolic valuation computation is fully retained, not replaced by this citation.
\item Earlier finite-field-parabolic argument and Bruhat counting. [Kos19] Automorphic Forms on the Stack of G-Zips, Results Math. 74 (2019), 91, DOI 10.1007/s00025-019-1021-z, Theorem 2, Theorem 3.5.3, Lemmas 3.4.4–3.4.5 and Proposition 3.5.2.. Author explicitly distinguishes the earlier finite-field-parabolic scope. Double-coset valuation invariance and current Frobenius orbit sum are reproduced through line 1125; finite Bruhat cardinalities remain cited tools.
\item Reductive-group realizations and induced representations. [Spr98] Sections 8.1.4 and proof of Lemma 9.2.2; [Jan03] Section 5.8; [IK24] Proposition 3.2.1.. Standard root SL2 maps, highest-weight induced modules and flag pushforward; definitions and the required identifications are retained in Section 2.
\end{itemize}\medskip
\ifdefined\hypertarget\hypertarget{sourcepointer-1}{}\fi
\noindent\textit{Source reference.} Original proposition 3.8: persistence-modules/\allowbreak{}EPIGA15861/\allowbreak{}sources/\allowbreak{}author-v5.tex, source lines 1012--1024. Forward-looking contextual reference in Section 2, not an input to the norm criterion; preserved as exact original-source pointer.\par\medskip
\noindent\textbf{Authors and source.} Naoki Imai and Jean-Stefan Koskivirta, \emph{Weights of mod $p$ automorphic forms and partial Hasse invariants}, exact arXiv:2211.16207v5 (12 July 2026), DOI \href{https://doi.org/10.46298/epiga.2026.15861}{10.46298/epiga.2026.15861}. Wushi Goldring authored the separate appendix, not the selected norm criterion.
\par\noindent\textbf{License.} Copyright the respective authors. Author text and this adapted presentation are under Creative Commons Attribution--ShareAlike 4.0 International, \url{https://creativecommons.org/licenses/by-sa/4.0/}. No endorsement is implied.
\par\noindent\textbf{Counted scope.} One result: entire Theorem 4 / Proposition 3.9, the highest-weight norm extension criterion for general cocharacter data, including parabolics not defined over the finite field. Introductory statement, complete Section 2, norm definition and all boundary/root-curve, Frobenius, valuation and Bruhat arguments through the end of the proof are retained. Supporting statements are not additional counted results.
\par\noindent\textbf{Changes and dependencies.} Journal furniture is replaced by a standard wrapper; original section, theorem and displayed-equation numbering is retained. The earlier results of Pink--Wedhorn--Ziegler, Imai--Koskivirta and Koskivirta are attributed external inputs, not represented as reproduced proofs. The current argument after the author's ``similar'' remark is reproduced in full. No mathematical repair or new proof is supplied.
\par\noindent\textbf{Preparation only.} No formal problem number or mathematical certification. Saved in-flight identities end on 14 October 2026 at 17:36:45 UTC; later overlap is unconfirmed. AI assistance concerns selection, metadata and packaging; no assertion about author AI use is inferred.

\setcounter{section}{1}\setcounter{thmIntro}{3}

\noeqref{eq:phiconj,orbparam,dualorbparam,VlambdadefH,Vlambda-function,Vlambda-natural-map,funct-rep,globquot,injection-ev1,Lphi-equ,eqL0,equ-Plowstar,equ-VDeltaP,zipflag-Cw,ident-H0-GF,eq:H0expid,Vflag-additivity,pullbackformula,chia-cond,zeta-Shimura,norm-def-eq,formula-norm,functionh-eq,eq-deltaalpha}

Our final result is the construction of natural mod $p$ automorphic forms attached to the highest-weight vectors of the representations $V_I(\lambda)$. Let $\lambda$ be an $L$-dominant character, and let $f_\lambda\in V_I(\lambda)$ denote the highest-weight vector of $V_I(\lambda)$. There is a natural way of defining the norm $\mathbf{f}_\lambda \colonequals \Norm_{L_\varphi}(f_\lambda)$ of $f_{\lambda}$. Here $L_\varphi$ is a certain finite (generally nonsmooth) subgroup of $L$ containing $L_0(\FF_q)$. There is an integer $m\geq 0$, determined by $L_\varphi$, such that the norm $\Norm_{L_\varphi}(f_\lambda)$ is a section of $\Vcal_I(d\lambda)$ (where $d=q^m|L_0(\FF_q)|$) over the $\mu$-ordinary locus $\Ucal_\mu$ of $\GZip^\mu$ (see Section~\ref{subsec-norm} for details). For $\alpha\in \Delta$, let $r_\alpha$ be the smallest integer $r\geq 1$ such that $\sigma^r(\alpha)=\alpha$.


\begin{thmIntro}[Proposition~\ref{prop-Norm}]\label{thm4}
The section $\mathbf{f}_\lambda$ extends to $\GZip^\mu$ if and only if for all $\alpha \in \Delta\setminus \Delta_L$ one has
\begin{equation}\label{formula-norm-intro}
\sum_{w\in W_{L_0}(\FF_q)} \sum_{i=0}^{r_\alpha-1} q^{i+\ell(w)} \ \left\langle w\lambda, \sigma^i\left(\alpha^\vee\right) \right\rangle\leq 0. \tag{1}
\end{equation}
\end{thmIntro}

\section{Preliminaries and background on the stack of \texorpdfstring{$\boldsymbol{G}$}{G}-zips}\label{sec2}

\subsection{Notation}\label{subsec-notation}

Throughout the paper, $p$ is a prime number, $q$ is a power of $p$, and $\FF_q$ is a finite field with $q$ elements. We write $k=\overline{\FF}_q$ for an algebraic closure of $\FF_q$. The notation $G$ will always denote a connected reductive group over $\FF_q$. For a $k$-scheme $X$, we denote by $X^{(q)}$ its $\supth{q}$ power Frobenius twist and by $\varphi \colon X\to X^{(q)}$ its relative Frobenius morphism. Write $\sigma \in \Gal(k/\FF_q)$ for the $q$-power Frobenius. We will always write $(B,T)$ for a Borel pair of $G$; \textit{i.e.}, $T \subset B \subset G$ are 
a maximal torus and a Borel subgroup in $G$. We do not assume that $T$ is split over $\FF_q$. Let $B^+$ be the Borel subgroup of $G$  opposite to $B$ with respect to $T$ (\textit{i.e.}, the unique Borel subgroup $B^+$ of $G$ such that $B^+\cap B=T$). We will use the following notation: 
\begin{itemize}
\item As usual, $X^*(T)$ (resp.~$X_*(T)$) denotes the group of characters (resp.~cocharacters) of $T$. The group $\Gal(k/\FF_q)$ acts naturally on these groups. Let $W=W(G_k,T)$ be the Weyl group of $G_k$. Similarly, $\Gal(k/\FF_q)$ acts on $W$. Furthermore, the actions of $\Gal(k/\FF_q)$ and $W$ on $X^*(T)$ and $X_*(T)$ are compatible in a natural sense. We write $W(\FF_q)$ for the $\Gal(k/\FF_q)$-fixed subgroup of $W$. 
\item We write $\Phi\subset X^*(T)$ for the set of $T$-roots of $G$.
\item We write $\Phi^+\subset \Phi$ for the system of positive roots with respect to $B^+$ (\textit{i.e.}, $\alpha \in \Phi^+$ when the $\alpha$-root group $U_{\alpha}$ is contained in $B^+$). This convention may differ from that of other authors. We use it to match the conventions of previous publications \cite{Goldring-Koskivirta-Strata-Hasse,Koskivirta-automforms-GZip}.
\item We write $\Delta\subset \Phi^+$ for the set of simple roots. 
\item For $\alpha \in \Phi$, let $s_\alpha \in W$ be the corresponding reflection. The system $(W,\{s_\alpha \mid \alpha \in \Delta\})$ is a Coxeter system. 
We write $\ell  \colon W\to \NN$ for the length function. Hence $\ell(s_\alpha)=1$ for all $\alpha\in \Delta$. Let $w_0$ denote the longest element of $W$.
\item For a subset $K\subset \Delta$, let $W_K$ denote the subgroup of $W$ generated by $\{s_\alpha \mid \alpha \in K\}$. Write $w_{0,K}$ for the longest element in $W_K$.
\item Let ${}^KW$ (resp.~$W^K$) denote the subset of elements $w\in W$ that have minimal length in the coset $W_K w$ (resp.~$wW_K$). Then ${}^K W$ (resp.\ $W^K$) is a set of representatives of $W_K\backslash W$ (resp.\ $W/W_K$). The map $g\mapsto g^{-1}$ induces a bijection ${}^K W\to W^K$. The longest element in the set ${}^K W$ is $w_{0,K} w_0$.
\item Let $X_{+}^*(T)$ denote the set of dominant characters, \textit{i.e.}, characters $\lambda\in X^*(T)$ such that $\langle \lambda,\alpha^\vee \rangle \geq 0$ for all $\alpha \in \Delta$.
\item For a subset $I\subset \Delta$, let $X_{+,I}^*(T)$ denote the set of characters $\lambda\in X^*(T)$ such that $\langle \lambda,\alpha^\vee \rangle \geq 0$ for all $\alpha \in I$. We call them $I$-dominant characters.
\item Let $P\subset G_k$ be a parabolic subgroup containing $B$, and let $L\subset P$ be the unique Levi subgroup of $P$ containing $T$. Then we define a subset $I_P\subset \Delta$ as the unique subset such that $W(L,T)=W_{I_P}$. For an arbitrary parabolic subgroup $P\subset G_k$ containing $T$, we define $I_P\subset \Delta$ as $I_P \colonequals I_{P'}$, where $P'$ is the unique conjugate of $P$ containing $B$.
\item For a parabolic $P\subset G_k$, write $\Delta^P \colonequals \Delta \setminus I_P$.
\item For all $\alpha\in \Phi$, choose an isomorphism $u_\alpha\colon \GG_{\mathrm{a}}\to U_\alpha$ so that 
  $(u_{\alpha})_{\alpha \in \Phi}$ is a realization in the sense of \cite[Section~8.1.4]{Springer-Linear-Algebraic-Groups-book}.
  In particular, we have 
\begin{equation}\label{eq:phiconj}
 t u_{\alpha}(x)t^{-1}=u_{\alpha}(\alpha(t)x), \quad \forall x\in \GG_{\mathrm{a}},\  \forall t\in T.
\end{equation}
\item Let $\phi_{\alpha}\colon \SL_2\to G$ denote the map attached to $\alpha$ as in \cite[Proof of Lemma~9.2.2]{Springer-Linear-Algebraic-Groups-book}.
  It satisfies
\[
 \phi_\alpha 
 \left( \begin{pmatrix}
 1 & x \\ 0 & 1 
 \end{pmatrix}\right) = u_{\alpha}(x), \quad 
 \phi_\alpha 
 \left( \begin{pmatrix}
 1 & 0 \\ x & 1 
 \end{pmatrix}\right) = u_{-\alpha}(x).
\]
\item Fix a $B$-representation $(V,\rho)$. For $j\in \ZZ$ and $\alpha\in \Phi$, we define a map $E_{\alpha}^{(j)} \colon V \to V$ as follows. Let $V=\bigoplus_{\nu \in X^*(T)}V_\nu$ be the weight decomposition of $V$. For $v\in V_\nu$, we can write uniquely
\[
 u_{\alpha}(x)v=\sum_{j \geq 0} x^j E_{\alpha}^{(j)}(v), \quad \forall x\in \GG_{\mathrm{a}},
\]
for elements $E_{\alpha}^{(j)}(v) \in V_{\nu+j\alpha}$ (see \cite[Lemma 3.3.1]{Imai-Koskivirta-vector-bundles}). Extend $E_{\alpha}^{(j)}$ by additivity to a map $V\to V$. For $j<0$, put $E_{\alpha}^{(j)}=0$. 
\end{itemize}

\subsection{The stack of $\boldsymbol{G}$-zips}\label{sec-Gzips}

We recall some facts about the stack of $G$-zips of Pink--Wedhorn--Ziegler in \cite{Pink-Wedhorn-Ziegler-zip-data}.

\subsubsection{Definitions} \label{subsec-zipdatum}
Let $G$ be a connected reductive group over $\FF_q$. In this paper, a zip datum is a tuple $\Zcal \colonequals (G,P,L,Q,M)$ consisting of the following objects:
\begin{equivlist}
    \item $P\subset G_k$ and $Q\subset G_k$ are parabolic subgroups of $G_k$.
    \item $L\subset P$ and $M\subset Q$ are Levi subgroups such that $L^{(q)}=M$. 
\end{equivlist}
For an algebraic group $H$, denote by $R_{\mathrm{u}}(H)$ the unipotent radical of $H$. If $P'\subset G_k$ is a parabolic subgroup with Levi subgroup $L'\subset P'$, any $x\in P'$ can be written uniquely as $x=\overline{x}u$ with $\overline{x}\in L'$ and $u\in R_{\mathrm{u}}(P')$. We denote by $\theta^{P'}_{L'} \colon P'\to L'$ the map $x\mapsto \overline{x}$. Since $M=L^{(q)}$, we have a Frobenius isogeny $\varphi \colon L\to M$. Put
\begin{equation}\label{zipgroup}
E \colonequals \left\{(x,y)\in P\times Q \relmiddle|  \varphi\left(\theta^P_L(x)\right)=\theta^Q_M(y)\right\}.
\end{equation}
Equivalently, $E$ is the subgroup of $P\times Q$ generated by $R_{\mathrm{u}}(P)\times R_{\mathrm{u}}(Q)$ and elements of the form $(a,\varphi(a))$ with $a\in L$. Let $G\times G$ act on $G$ by $(a,b)\cdot g \colonequals agb^{-1}$, and let $E$ act on $G$ by restricting this action to $E$. The stack of $G$-zips of type $\Zcal$ (see \cite{Pink-Wedhorn-Ziegler-zip-data,Pink-Wedhorn-Ziegler-F-Zips-additional-structure}) can be defined as the quotient stack
\[
\GZip^\Zcal = \left[E\backslash G_k \right].
\]

\subsubsection{Cocharacter datum} \label{subsec-cochar}
A \emph{cocharacter datum} is a pair $(G,\mu)$, where $G$ is a reductive connected group over $\FF_q$ and $\mu \colon \GG_{\mathrm{m},k}\to G_k$ is a cocharacter. One can attach to $(G,\mu)$ a zip datum $\Zcal_\mu$, defined as follows. First, denote by $P_+(\mu)$ (resp.\ $P_-(\mu)$) the unique parabolic subgroup of $G_k$ such that $P_+(\mu)(k)$ (resp.\ $P_-(\mu)(k)$) consists of the elements $g\in G(k)$ satisfying that the map 
\[
\GG_{\mathrm{m},k} \lra G_{k},\quad  t\longmapsto\mu(t)g\mu(t)^{-1} \quad \left(\textrm{resp.\ } t\longmapsto\mu(t)^{-1}g\mu(t)\right)
\]
extends to a morphism of varieties $\AA_{k}^1\to G_{k}$. We obtain a pair of parabolics $(P_+(\mu),P_{-}(\mu))$ in $G_k$ whose intersection $P_+(\mu)\cap P_-(\mu)=L(\mu)$ is the centralizer of $\mu$ (it is a common Levi subgroup of $P_+(\mu)$ and $P_-(\mu)$). Set $P \colonequals P_-(\mu)$, $Q \colonequals (P_+(\mu))^{(q)}$, $L \colonequals L(\mu)$ and $M \colonequals  (L(\mu))^{(q)}$. The tuple $\Zcal_\mu \colonequals (G,P,L,Q,M)$ is a zip datum, which we call the zip datum attached to the cocharacter datum $(G,\mu)$. We write simply $\GZip^\mu$ for $\GZip^{\Zcal_\mu}$. We always consider zip data of this form.

\begin{rmk}\label{rmk-opposite-cochar}
A general zip datum $(G,P,L,Q,M)$ is of the form $\Zcal_\mu$ for a cocharacter $\mu \colon \GG_{\mathrm{m},k}\to G_k$ if and only if $\sigma(P)$ and $Q$ are opposite parabolic subgroups with common Levi $M=\sigma(L)$.
\end{rmk}

\begin{rmk}\label{rmkmuFp}
If $\mu$ is defined over $\FF_q$, then so are $P$ and $Q$. In this case, we have $L=M$, and $P$ and $Q$ are opposite parabolic subgroups with common Levi subgroup $L$.
\end{rmk}

\subsubsection{Frames} \label{sec-frames}
Let $\Zcal=(G,P,Q,L,M)$ be a zip datum. In this paper, a frame for $\Zcal$ is a triple $(B,T,z)$, where $(B,T)$ is a Borel pair of $G_k$ defined over $\FF_q$ such that 
\begin{equivlist}
    \item one has the inclusion $B\subset P$; 
    \item $z\in W$ is an element satisfying the conditions
\begin{equation}\label{eqBorel}
{}^z \! B \subset Q \quad \textrm{and} \quad
 B\cap M= {}^z \! B\cap M. 
\end{equation}
\end{equivlist}
We put $B_M\colonequals B\cap M$. Other papers (see \cite{Pink-Wedhorn-Ziegler-zip-data,Pink-Wedhorn-Ziegler-F-Zips-additional-structure, Koskivirta-Wedhorn-Hasse}) use the convention $B\subset Q$ instead of $B\subset P$. A frame (as defined here) may not always exist. However, if $(G,\mu)$ is a cocharacter datum and $\Zcal_\mu$ is the associated zip datum by Section~\ref{subsec-cochar}, then there exists a $G(k)$-conjugate $\mu'=\ad(g)\circ \mu$ (with $g\in G(k)$) such that $\Zcal_{\mu'}$ admits a frame, by Lemma~\ref{equ-IJDeltaP} below. Hence, it is harmless to assume that a frame exists, and we only consider zip data that admit frames. With respect to the Borel pair $(B,T)$, we define subsets $I,J,\Delta^P $ of~$\Delta$ as follows:
\begin{equation}\label{equ-IJDeltaP}
 I \colonequals I_P, \quad J \colonequals I_Q, \quad \Delta^P=\Delta\setminus I.  
\end{equation}

\begin{lemma}[\textit{cf.} {\cite[Lemma 2.3.4]{Goldring-Koskivirta-zip-flags}}]\label{lem-framemu}
Let $\mu \colon \GG_{\mathrm{m},k}\to G_k$ be a cocharacter, and let $\Zcal_\mu$ be the attached zip datum. Assume that $(B,T)$ is a Borel pair defined over $\FF_q$ such that $B\subset P$. Define the element
\begin{equation}\label{z-def}
z \colonequals w_0 w_{0,J}=\sigma(w_{0,I})w_0.
\end{equation}
Then $(B,T,z)$ is a frame for $\Zcal_\mu$.
\end{lemma}

\subsubsection{Parametrization of the \texorpdfstring{$\boldsymbol{E}$}{E}-orbits in \texorpdfstring{$\boldsymbol{G}$}{G}}
\label{subsec-zipstrata}
By \cite[Proposition 7.1]{Pink-Wedhorn-Ziegler-zip-data}, there are finitely many $E$-orbits in $G$. The $E$-orbits are smooth and locally closed in $G$, and the Zariski closure of an $E$-orbit is a union of $E$-orbits. We review the parametrization of $E$-orbits following \cite{Pink-Wedhorn-Ziegler-zip-data}. For $w\in W$, fix a representative $\dot{w}\in N_G(T)$ such that $(w_1w_2)^\cdot = \dot{w}_1\dot{w}_2$ whenever $\ell(w_1 w_2)=\ell(w_1)+\ell(w_2)$ (this is possible by choosing a Chevalley system; see \cite[Section~XXIII.6]{SGA3}). For $w\in W$, define $G_w$ as the $E$-orbit of $\dot{w}\dot{z}^{-1}$. If no confusion occurs, we write $w$ instead of $\dot{w}$. For $w,w'\in {}^I W$, write $w'\preccurlyeq w$ if there exists a $w_1\in W_I$ such that $w'\leq w_1 w \sigma(w_1)^{-1}$. This defines a partial order on ${}^I W$ (see \cite[Corollary 6.3]{Pink-Wedhorn-Ziegler-zip-data}).

\begin{theorem}[\textit{cf.} {\cite[Theorems 7.5,~11.2,~11.3 and~11.5]{Pink-Wedhorn-Ziegler-zip-data}}] \label{thm-E-orb-param}
We have two bijections:
\begin{align} \label{orbparam}
{}^I W &\longrightarrow \{ \textrm{$E$-orbits in $G_k$} \}, \quad w\longmapsto G_w,  \\  
\label{dualorbparam} W^J &\longrightarrow \{ \textrm{$E$-orbits in $G_k$}\}, \quad  w\longmapsto G_w. 
\end{align}
For $w\in {}^I W\cup W^J$, one has $\dim(G_w)= \ell(w)+\dim(P)$, and the Zariski closure of $G_w$ is 
\begin{equation}\label{equ-closure-rel}
\overline{G}_w=\bigsqcup_{w'\in {}^IW,\  w'\preccurlyeq w} G_{w'}
\end{equation}
for $w\in {}^I W$ and 
\begin{equation}\label{equ-closure-relJ}
\overline{G}_w=\bigsqcup_{w'\in W^J,\  w'\preccurlyeq w} G_{w'}
\end{equation}
for $w\in W^J$.
\end{theorem}
In particular, there is a unique open $E$-orbit $U_\Zcal\subset G$ corresponding to the longest elements $w_{0,I}w_0\in {}^I W$ via \eqref{orbparam} and to $w_0w_{0,J}\in W^J$ via \eqref{dualorbparam}. The $E$-orbit $U_\Zcal$ is dense in $G$.  If $\Zcal=\Zcal_\mu$ (see Section~\ref{subsec-cochar}), write $U_\mu=U_{\Zcal_\mu}$. In this case, we can choose $z=w_{0}w_{0,J}=\sigma(w_{0,I})w_0$ (see Lemma~\ref{lem-framemu}); hence \eqref{dualorbparam} shows that $1\in U_\mu$.  We put $\Ucal_\mu \colonequals [E\backslash U_\mu]$, which we call the $\mu$-ordinary locus.

\subsection{Vector bundles on the stack of $\boldsymbol{G}$} \label{sec-vector-bundles-gzipz}

\subsubsection{Representation theory}\label{subsec-remind}

For an algebraic group $G$ over a field $K$, denote by $\Rep(G)$ the category of algebraic representations of $G$ on finite-dimensional $K$-vector spaces. We denote a representation $\rho\colon G\to \GL_K(V)$ by $(V,\rho)$, or sometimes simply $\rho$ or $V$. For an algebraic group $G$ over $\FF_q$, a $G_k$-representation $(V,\rho)$ and an integer $m$, we denote by $(V^{[m]},\rho^{[m]})$ the representation such that $V^{[m]}=V$ and
\begin{equation}\label{equ-rhom}
  \rho^{[m]} \colon  G_k \xrightarrow{\varphi^m} G_k \overset{\rho}\lra \GL(V).
\end{equation}

Let $H$ be a split connected reductive $K$-group, and choose a Borel pair $(B_H,T)$ defined over $K$. If $K$ has characteristic zero, $\Rep(H)$ is semisimple. In characteristic $p$, however, this is no longer true in general. For $\lambda\in X_{+}^*(T)$, let $\Lcal_\lambda$ be the line bundle attached to $\lambda$ on the flag variety $H/B_H$ by the usual associated sheaf construction (see \cite[Section~5.8]{jantzen-representations}). Define an $H$-representation $V_H(\lambda)$ by
\begin{equation}\label{VlambdadefH}
    V_H(\lambda) \colonequals H^0\left(H/B_H,\Lcal_\lambda\right) . 
\end{equation}
In other words, one has $V_H(\lambda)=\Ind_{B_H}^{H} \lambda$. The representation $V_H(\lambda)$ is of highest weight $\lambda$. If $\cara(K)=0$, the representation $V_H(\lambda)$ is irreducible. We view elements of $V_H(\lambda)$ as regular maps $f \colon H\to \AA^1$ satisfying
\begin{equation}\label{Vlambda-function}
f(hb)=\lambda\left(b^{-1}\right)f(h), \quad \forall h\in H, \ \forall b\in B_H.    
\end{equation}
For dominant characters $\lambda,\lambda'$, there is a natural surjective map
\begin{equation}\label{Vlambda-natural-map}
V_H(\lambda)\otimes V_H(\lambda')\lra V_H(\lambda+\lambda').
\end{equation}
In the description given by \eqref{Vlambda-function}, this map is $f\otimes f'\mapsto ff'$ (for $f\in V_H(\lambda)$, $f'\in V_H(\lambda')$). Denote by $W_H \colonequals W(H,T)$ the Weyl group and by $w_{0,H}\in W_H$ the longest element. Then $V_H(\lambda)$ has a unique $B_H$-stable line, which is a weight space for the weight $w_{0,H}\lambda$. 


\subsubsection{Vector bundles on quotient stacks} \label{sec-genth-vb}
For an algebraic stack $\Xcal$, write $\VB(\Xcal)$ for the category of vector bundles on $\Xcal$. Let $X$ be a $k$-scheme and $H$ an affine $k$-group scheme acting on $X$. If $\rho \colon H\to \GL(V)$ is an algebraic representation of $H$, it gives rise to a vector bundle $\Vcal_{H,X}(\rho)$ on the stack $[H\backslash X]$. This vector bundle can be defined geometrically as $[H\backslash (X\times_k V)]$, where $H$ acts diagonally on $X\times_k V$. We obtain a functor
\begin{equation}\label{funct-rep}
    \Vcal_{H,X} \colon \Rep(H)\lra \VB([H\backslash X]).
\end{equation}
Similarly to the usual associated sheaf construction, see \cite[Section~5.8, Equation (1)]{jantzen-representations}, the global sections of $\Vcal_{H,X}(\rho))$ are given by
\begin{equation}\label{globquot}
 H^0([H\backslash X],\Vcal_{H,X}(\rho))
 =\{f \colon X\to V \mid  f(h \cdot x)=\rho(h) f(x) , \  \forall h\in H, \ \forall x\in X\},
\end{equation}
where $f\colon X\to V$ is a morphism of $k$-schemes, and $V$ is viewed as an affine space over $k$.

\subsubsection{Vector bundles on \texorpdfstring{$\boldsymbol{\GZip^\mu}$}{GZip\textasciicircum\{mu\}}} \label{sec-VB-Gzip}

Fix a cocharacter datum $(G,\mu)$. Let $\Zcal=(G,P,L,Q,M)$ be the attached zip datum. Fix a frame $(B,T)$ as in Section~\ref{sec-frames}. By \eqref{funct-rep}, we have a functor $\Vcal_{E,G}\colon \Rep(E)\to \VB(\GZip^\mu)$, which we simply denote by $\Vcal$. For $(V,\rho)\in \Rep(E)$, the global sections of $\Vcal(\rho)$ are
\[
 H^0(\GZip^\mu,\Vcal(\rho))=\left\{f \colon G_k \to V \mid f(\epsilon \cdot g)=\rho(\epsilon) f(g) , \ \forall\epsilon\in E, \ \forall g\in G_k \right\}.
\]
Since $G$ admits an open dense $E$-orbit (see the discussion below Theorem~\ref{thm-E-orb-param}), the space $H^0(\GZip^\mu,\Vcal(\rho))$ is finite-dimensional (see \cite[Lemma 1.2.1]{Koskivirta-automforms-GZip}). The first projection $p_1 \colon E\to P$ induces a functor $p_1^* \colon \Rep(P)\to \Rep(E)$. If $(V,\rho)\in \Rep(P)$, we write again $\Vcal(\rho)$ for $\Vcal(p_1^*(\rho))$. In this paper, we only consider $E$-representations coming from $P$ in this way. Let $\theta^P_L \colon P\to L$ be the natural projection modulo $R_{\mathrm{u}}(P)$, as in Section~\ref{subsec-zipdatum}. It induces a fully faithful functor
\begin{equation}
(\theta^P_L)^* \colon \Rep(L)\lra \Rep(P)  
\end{equation}
whose image is the full subcategory of $\Rep(P)$ of $P$-representations trivial on $R_{\mathrm{u}}(P)$. Hence, we view $\Rep(L)$ as a full subcategory of $\Rep(P)$. If $(V,\rho)\in \Rep(L)$, write again $\Vcal(\rho) \colonequals \Vcal((\theta^P_L)^*\rho)$. For $\lambda\in X^*(T)$, write $B_L \colonequals B\cap L$, and define an $L$-representation $(V_I(\lambda),\rho_{I,\lambda})$ as follows: 
\begin{equation}\label{equ-VlambdaL}
V_I(\lambda)=\Ind_{B_L}^L \lambda, \quad \rho_{I,\lambda}\colon L\lra \GL(V_I(\lambda)). 
\end{equation}
This is the representation defined in \eqref{VlambdadefH} for $H=L$ and $B_H=B_L$. Let $\Vcal_I(\lambda)$ be the vector bundle on $\GZip^\mu$ attached to $V_I(\lambda)$, and call it an \emph{automorphic vector bundle} on $\GZip^\mu$ associated to $\lambda$. This terminology stems from Shimura varieties (see Section~\ref{subsec-Shimura} below for further details). For $\lambda \in X^*(L)$, viewing $\lambda$ as an element of $X^*(T)$ by restriction, the vector bundle $\Vcal_I(\lambda)$ is a line bundle. Note that if $\lambda\in X^*(T)$ is not $I$-dominant, then $V_I(\lambda)=0$ and thus $\Vcal_I(\lambda)=0$.

\subsection{Global sections over $\boldsymbol{\GZip^\mu}$} \label{subsec-global-sections}
We review some results of \cite{Imai-Koskivirta-vector-bundles} regarding the global sections of $\Vcal(\rho)$ for a $P$-representation $\rho$. We start with sections over the open substack $\Ucal_\mu\subset \GZip^\mu$. Recall that $\Ucal_\mu=[E\backslash U_\mu]$ and $1\in U_\mu$ (see Section~\ref{subsec-zipstrata}). By \eqref{globquot}, an element of $H^0(\Ucal_\mu,\Vcal(\rho))$ can be viewed as a map $h\colon G\to V$ satisfying $h(agb^{-1})=\rho(a)h(g)$ for all $(a,b)\in E$ and all $g\in G$. Since the $E$-orbit of $1$ is open dense in $G$, the map $h\mapsto h(1)$ is an injection
\begin{equation}\label{injection-ev1}
\ev_1 \colon H^0\left(\Ucal_\mu,\Vcal(\rho)\right)\lra V.
\end{equation}
We give the image of this map. Let $L_\varphi$ be the scheme-theoretic stabilizer subgroup of $1$ in $E$. By definition, one has
\begin{equation}\label{Lphi-equ}
L_{\varphi}=E\cap \{(x,x) \mid x\in G_k \}, 
\end{equation}
which is a $0$-dimensional algebraic group (in general nonsmooth). The first projection $E\to P$ induces a closed immersion $L_{\varphi}\to P$. Identify $L_\varphi$ with its image, and view it as a subgroup of $P$. Denote by $L_0\subset L$ the largest algebraic subgroup defined over $\FF_q$. 
In other words,
\begin{equation}\label{eqL0}
    L_0=\bigcap_{n\geq 0}L^{(q^n)}.
\end{equation}

\begin{lemma}[\textit{cf.} {\cite[Lemma 3.2.1]{Koskivirta-Wedhorn-Hasse}}]\label{lemLphi} \leavevmode
\begin{assertionlist}
\item \label{lemLphi-item1} One has $L_{\varphi}\subset L$.
\item \label{lemLphi-item2}  The group $L_{\varphi}$ can be written as a semidirect product $L_{\varphi}=L_{\varphi}^\circ\rtimes L_0(\FF_q)$, where $L_{\varphi}^\circ$ is the identity component of\, $L_{\varphi}$. Furthermore, $L_{\varphi}^\circ$ is a finite unipotent algebraic group.
\item \label{lemLphi-item3}  Assume that $P$ is defined over $\FF_q$. Then $L_0=L$ and $L_{\varphi}=L(\FF_q)$, viewed as a constant algebraic group.
\end{assertionlist}
\end{lemma}

\begin{lemma}[\textit{cf.} {\cite[Corollary 3.2.3]{Imai-Koskivirta-vector-bundles}}] \label{lem-Umu-sections}
The map \eqref{injection-ev1} induces an identification
\begin{equation}
 H^0\left(\Ucal_\mu,\Vcal(\rho)\right)=V^{L_\varphi}.   
\end{equation}
\end{lemma}

Here, the notation $V^{L_{\varphi}}$ denotes the space of scheme-theoretic invariants, \textit{i.e.}, the set of $v\in V$ such that for any $k$-algebra $R$, one has $\rho(x)v=v$ in $V\otimes_k R$ for all $x\in L_{\varphi}(R)$. We now consider the space of global sections over $\GZip^\mu$. Restriction of sections to $\Ucal_\mu\subset \GZip^\mu$ induces an injective map $H^0(\GZip^{\mu},\Vcal(\rho))\to H^0(\Ucal_\mu,\Vcal(\rho))=V^{L_\varphi}$. For simplicity, we assume here that $P$ is defined over $\FF_q$ (for the general result, see \cite[Theorem 3.4.1]{Imai-Koskivirta-vector-bundles}). We will need the general version in the proof of Proposition~\ref{prop-charL}, but in the simple setting when $\rho$ is a character $L\to \GG_{\mathrm{m}}$. For $\alpha\in \Phi$, choose a realization $(u_{\alpha})_{\alpha \in \Phi}$ (see Section~\ref{subsec-notation}). Fix a $P$-representation $(V,\rho)$, and let $V=\bigoplus_{\nu \in X^*(T)} V_\nu$ be its $T$-weight decomposition. Define the Brylinski--Kostant filtration (\cf \cite[Equation~(3.3.2)]{Xiao-Zhu-on-vector-valued}) indexed by $c\in \RR$ on $V_\nu$ by 
\begin{equation}\label{equ-Filc}
\fil_{c}^{\alpha} V_{\nu} = 
 \bigcap_{j > c} \Ker \left( 
 E_{\alpha}^{(j)} \colon V_{\nu} \to V_{\nu+j\alpha} \right), 
\end{equation}
where the map $E_{\alpha}$ was defined in Section~\ref{subsec-notation}. For $\chi \in X^*(T)_{\RR}$ and $\nu \in X^*(T)$, also set
\begin{equation}\label{equ-FilPchi}
 \fil_{\chi}^P 
 V_{\nu} = 
 \bigcap_{\alpha \in \Delta^P} 
 \fil_{\langle \chi, \alpha^\vee \rangle}^{-\alpha} 
 V_{\nu}.
\end{equation}
The Lang torsor morphism $\wp \colon T \to T$, $g\mapsto g\varphi(g)^{-1}$ induces isomorphisms 
\begin{align}
 &\wp^* \colon X^*(T)_{\RR} \stackrel{\lowsim}{\longrightarrow} X^*(T)_{\RR},\quad \lambda \longmapsto \lambda \circ \wp  = \lambda -q\sigma^{-1}(\lambda), \label{equ-Pupstar} \\
 &\wp_* \colon X_*(T)_{\RR} \stackrel{\lowsim}{\longrightarrow} X_*(T)_{\RR},\quad \delta \longmapsto \wp \circ \delta  = \delta -q\sigma(\delta).  \label{equ-Plowstar}
\end{align}

\begin{theorem}[\textit{cf.} {\cite[Corollary 3.4.2]{Imai-Koskivirta-vector-bundles}}]\label{main-Fq}
Assume that $P$ is defined over $\FF_q$. 
For all $(V,\rho)\in \Rep(P)$, the map $\ev_1$ induces an identification
\begin{equation}
H^0(\GZip^\mu,\Vcal(\rho))=V^{L(\FF_q)}\cap 
 \bigoplus_{\nu \in X^*(T)} 
 \fil_{{\wp^*}^{-1}(\nu)}^P 
 V_{\nu}.
\end{equation}
\end{theorem}

In the general case of an arbitrary parabolic $P$, $V^{L(\FF_q)}$ is replaced by $V^{L_\varphi}$ and $\Fil^\alpha_c V_\nu$ is replaced by a generalized Brylinski--Kostant filtration (see \cite[Theorem 3.4.1]{Imai-Koskivirta-vector-bundles}). In the special case when $\rho$ is trivial on $R_{\mathrm{u}}(P)$, Theorem~\ref{main-Fq} simplifies greatly. Set $\delta_{\alpha}\colonequals \wp_*^{-1}(\alpha^\vee)$, and define a subspace $V_{\geq 0}^{\Delta^P}\subset V$ by

\begin{equation}\label{equ-VDeltaP}
V_{\geq 0}^{\Delta^P} = \bigoplus_{\substack{\langle \nu,\delta_{\alpha} \rangle \geq 0, \ 
 \forall \alpha\in \Delta^P}} V_\nu.
\end{equation}
If $T$ is split over $\FF_q$, then 
$\delta_{\alpha}=-\alpha^\vee /(q-1)$, 
and $V_{\geq 0}^{\Delta^P}$ is the direct sum of the weight spaces $V_\nu$ for those $\nu\in X^*(T)$ satisfying $\langle \nu,\alpha^\vee \rangle \leq 0$ for all $\alpha \in \Delta^P$. 


\begin{corollary}\label{cor-Fq-Levi}
Assume that $P$ is defined over $\FF_q$ and furthermore that $(V,\rho)\in \Rep(P)$ is trivial on $R_{\mathrm{u}}(P)$. Then one has
\begin{equation}
H^0(\GZip^\mu,\Vcal(\rho))=V^{L(\FF_q)}\cap V_{\geq 0}^{\Delta^P} .
\end{equation}
\end{corollary}



\subsection{The stack of $\boldsymbol{G}$-zip flags} \label{subsec-zipflag}

\subsubsection{Definition}
Let $(G,\mu)$ be a cocharacter datum with attached zip datum $\Zcal_{\mu}=(G,P,L,Q,M)$ (see Section~\ref{subsec-cochar}). Fix a frame $(B,T,z)$ with $z=\sigma(w_{0,I})w_0=w_0w_{0,J}$ (see Lemma~\ref{lem-framemu}). The stack of zip flags (see \cite[Definition 2.1.1]{Goldring-Koskivirta-Strata-Hasse}) is defined as
\begin{equation}\label{eq-Gzipflag-PmodB}
\GF^\mu=[E\backslash (G_k \times P/B)],     
\end{equation}
where the group $E$ acts on the variety $G_k \times (P/B)$ by the rule $(a,b)\cdot (g,hB) \colonequals (agb^{-1},ahB)$ for all $(a,b)\in E$ and all $(g,hB)\in G_k \times P/B$. The first projection $G_k \times P/B \to G_k$ is $E$-equivariant and yields a natural morphism of stacks
 \begin{equation}\label{projmap-flag}
     \pi  \colon  \GF^\mu \lra \GZip^\mu.
 \end{equation}

Set $E' \colonequals E\cap (B\times G_k)$. Then the injective map $G_k \to G_k \times P/B$, $g\mapsto (g,B)$ yields an isomorphism of stacks $[E' \backslash G_k]\simeq \GF^\mu$ (see \cite[Equation~(2.1.5)]{Goldring-Koskivirta-Strata-Hasse}). We recall the stratification of $\GF^\mu$. First, define the Schubert stack as the quotient stack
\begin{equation}\label{equ-def-Sbt}
\Sbt \colonequals [B\backslash G_k /B].
\end{equation}
This stack is finite and smooth. Its topological space is isomorphic to $W$, endowed with the topology induced by the Bruhat order on $W$. This follows easily from the Bruhat decomposition of $G$. One can show that $E'\subset B\times {}^z B$. In particular, there is a natural projection map $[E'\backslash G_k]\to [B\backslash G_k/{}^z\! B]$. Composing with the isomorphism $[B\backslash G_k/{}^z B]\to [B\backslash G_k/B]$ induced by $G_k\to G_k$, $g\mapsto gz$, we obtain a smooth, surjective map
 \begin{equation}\label{eq-GF-to-Sbt}
     \psi  \colon  \GF^\mu \lra \Sbt.
 \end{equation}
For $w\in W$, put $\Sbt_w\colonequals [B\backslash BwB /B]$; it is a locally closed substack of $\Sbt$. The \emph{flag strata} of $\GF^\mu$ are defined as fibers of the map $\psi$. They are locally closed substacks (endowed with the reduced structure). Concretely, let $w\in W$, and put
\begin{equation}
F_w \colonequals B\left(wz^{-1}\right){}^zB=BwBz^{-1},
\end{equation} 
which is the $B\times {}^zB$-orbit of $wz^{-1}$. The set $F_w$ is locally closed in $G_k$, and one has $\dim(F_w)=\ell(w)+\dim(B)$. Then, via the isomorphism $\GF^\mu\simeq [E'\backslash G_k]$, the flag strata of $\GF^\mu$ are the locally closed substacks
\begin{equation}\label{zipflag-Cw}
\Fcal_w \colonequals [E'\backslash F_w], \quad w\in W.    
\end{equation}
The set $F_{w_0}\subset G_k$ is open in $G_k$, and similarly the stratum $\Fcal_{w_0}$ is open in $\GF^\mu$. The $B\times {}^zB$-orbits of codimension $1$ are $F_{s_\alpha w_0}$ for $\alpha\in \Delta$. The Zariski closure $\overline{F}_w$ is normal (see \cite[Theorem 3]{Ramanan-Ramanathan-projective-normality}) and coincides with $\bigcup_{w'\leq w}F_{w'}$, where $\leq$ is the Bruhat order of $W$.

\subsubsection{Vector bundles on  \texorpdfstring{$\boldsymbol{\GF^\mu}$}{GF\textasciicircum mu}}
Let $\rho \colon B\to \GL(V)$ be an algebraic representation, and view~$\rho$ as a representation of $E'$ via the first projection $E'\to B$. Via the isomorphism $\GF^\mu\simeq [E'\backslash G_k]$, we obtain a vector bundle $\Vcal_{\rm flag}$ on $\GF^\mu$. Let $(V,\rho)\in \Rep(P)$, and let $\Vcal(\rho)$ be the attached vector bundle on $\GZip^\mu$. Then one has
\begin{equation}\label{pullbackV}
\pi^*(\Vcal(\rho))=\Vcal_{\flag}(\rho|_B).
\end{equation}
Note that the rank of $\Vcal_{\flag}(\rho)$ is the dimension of $\rho$. In particular, if $\lambda\in X^*(B)$, then $\Vcal_{\flag}(\lambda)$ is a line bundle. For $(V,\rho)\in \Rep(B)$, consider the $P$-representation $\Ind_B^P(\rho)$ defined by
\begin{equation}\label{eq-Indrho}
\Ind_B^P(\rho)=\left\{f\colon P\to V \relmiddle| f(xb) = \rho\left(b^{-1}\right)f(x), \ \forall b\in B,\ \forall x\in P \right\}.
\end{equation}
For $y\in P$ and $f\in \Ind_B^P(\rho)$, the element $y\cdot f$ is the function $x\mapsto f(y^{-1} x)$.

\begin{proposition}[\textit{cf.} {\cite[Proposition 3.2.1]{Imai-Koskivirta-partial-Hasse}}]\label{prop-push-flag-zip}
For $(V,\rho)\in \Rep(B)$, we have the identification 
$\pi_{*}(\Vcal_{\flag}(\rho))=\Vcal(\Ind_B^P(\rho))$. In particular $\pi_{*}(\Vcal_{\flag}(\rho))$ is a vector bundle on $\GZip^\mu$.
\end{proposition}

In particular, if $\rho$ is a character $\lambda\in X^*(T)$, then $\Vcal_{\flag}(\lambda)$ is a line bundle and one has 
\begin{equation}\label{pushforward-lambda}
    \pi_*\left(\Vcal_{\flag}(\lambda)\right)=\Vcal_I(\lambda), 
\end{equation}
where the vector bundle $\Vcal_I(\lambda)$ was defined in Section~\ref{sec-VB-Gzip}. Hence, we have 
\begin{equation}\label{ident-H0-GF}
    H^0(\GZip^\mu,\Vcal_I(\lambda))=H^0\left(\GF^\mu,\Vcal_{\flag}(\lambda)\right).
\end{equation}
If $f \colon G_k \to k$ is a section of the right-hand side of \eqref{ident-H0-GF}, 
then the corresponding function 
$f_I \colon G_k \to V_I (\lambda)$ 
on the left-hand side of 
\eqref{ident-H0-GF} is given by 
\begin{equation}\label{eq:H0expid}
 (f_I (g))(x)=f \left(\left(x^{-1},\varphi(x)^{-1}\right) \cdot g\right)= 
 f\left(x^{-1} g \varphi (x)\right) 
\end{equation}
for all $g \in G_k$ and $x \in L$, by the construction of the identification. 
Also note  that the line bundles $\Vcal_{\flag}(\lambda)$ satisfy the following identity:
\begin{equation}\label{Vflag-additivity}
\Vcal_{\flag}(\lambda+\lambda')=\Vcal_{\flag}(\lambda)\otimes \Vcal_{\flag}(\lambda'), \quad \forall \lambda, \lambda'\in X^*(T).
\end{equation}
We can also define vector bundles on the stack $\Sbt$ as in \cite[Section~4]{Imai-Koskivirta-partial-Hasse}. For our purpose, it is enough to define line bundles on $\Sbt$. Using \eqref{funct-rep}, we can attach to each $(\chi_1,\chi_2)\in X^*(T)\times X^*(T)$ a line bundle $\Vcal_{\Sbt}(\chi_1,\chi_2)$ on $\Sbt$. One has
\begin{equation}\label{pullbackformula}
\psi^*\Vcal_{\Sbt}(\chi_1,\chi_2)=\Vcal_{\flag}(\chi_1+(z\chi_2)\circ \varphi)=\Vcal_{\flag}\left(\chi_1+q\sigma^{-1}(z\chi_2)\right).
\end{equation}

\subsubsection{Partial Hasse invariants}\label{sec-partialHasse}
We recall some results of \cite{Imai-Koskivirta-partial-Hasse}. By \cite[Theorem 2.2.1(a)]{Goldring-Koskivirta-Strata-Hasse}, the line bundle $\Vcal_{\Sbt}(\chi_1,\chi_2)$ admits a nonzero section over $\Sbt_{w_0}$ if and only if $\chi_1=-w_0\chi_2$. If this condition is satisfied, $H^0(\Sbt_{w_0}, \Vcal_{\Sbt}(\chi_1,\chi_2))$ is $1$-dimensional. For $\chi\in X^*(T)$, let $h_\chi$ be any nonzero element
\begin{equation}\label{hchi}
  h_\chi\in H^0\left(\Sbt_{w_0}, \Vcal_{\Sbt}(-w_0\chi,\chi)\right).
\end{equation}
By \cite[Theorem 2.2.1(c)]{Goldring-Koskivirta-Strata-Hasse}, $h_\chi$ extends to $\Sbt$ if and only if $\chi$ is a dominant character. Using \eqref{pullbackformula} and $z=\sigma(w_{0,I})w_0$, we obtain a section
\begin{equation}
\Ha_{\chi}\colonequals \psi^*(h_\chi)\in H^0\left(\Fcal_{w_0},\Vcal_{\flag}\left(-w_0\chi+qw_{0,I}w_0\left(\sigma^{-1}\chi\right)\right)\right), 
\end{equation}
and for $\chi\in X^*_+(T)$ the section $\Ha_{\chi}$ extends to $\GF^\mu$. In particular, let $\alpha\in \Delta$, and suppose $\chi_{\alpha}$ is a character satisfying 
\begin{equation}\label{chia-cond}
\begin{cases}
\left\langle \chi_\alpha, \alpha^{\vee} \right\rangle>0, & \\
\left\langle \chi_\alpha, \beta^{\vee} \right\rangle=0 & \ \textrm{ for all } \beta\in \Delta\setminus \{\alpha\}.
\end{cases}
\end{equation}
In this case, the section $h_{\chi_\alpha}$ vanishes exactly on the codimension~$1$ stratum $\overline{\Sbt}_{w_0s_\alpha}$. Similarly, the section $\Ha_{\chi_{\alpha}}$ cuts out the Zariski closure of the codimension~$1$ stratum $\Fcal_{w_0s_\alpha}$.


\begin{definition}\label{partial-Hasse-def}
For $\alpha\in \Delta$ and $\chi_\alpha$ satisfying \eqref{chia-cond}, we call the section $\Ha_{\chi_{\alpha}}$ a partial Hasse invariant for the stratum $\Fcal_{w_0 s_\alpha}$.
\end{definition}

\subsection{Shimura varieties and $\boldsymbol{G}$-zips} \label{subsec-Shimura}
We explain the connection between the stack of $G$-zips and Shimura varieties. Let $\gx$ be a Shimura datum; see \cite[Section~2.1.1]{Deligne-Shimura-varieties}. 
Write $\mathbf{E}=\egx$ for the reflex field of $\gx$ and $\Ocal_\mathbf{E}$ for its ring of integers. 
Given an open compact subgroup $K \subset \gofaf$, write $\shgx_{K}$ for Deligne's canonical model at level $K$ over $\mathbf{E}$ (see \cite{Deligne-Shimura-varieties}). For $K\subset \mathbf{G}(\AA_f)$ small enough, $\shgx_K$ is a smooth, quasi-projective scheme over $\mathbf{E}$. Every inclusion $K' \subset K$ induces a finite \'{e}tale projection $\pi_{K'/K} \colon \shgx_{K'} \to \shgx_{K}$.


Fix a prime number $p$, and assume that the level $K$ is of the form $K=K_pK^p$,  where $K_p\subset \mathbf{G}(\QQ_p)$ is a hyperspecial subgroup and $K^p\subset \mathbf{G}(\AA_f^p)$ is an open compact subgroup. Then one has $K_p=\Gscr(\ZZ_p)$, where $\Gscr$ is a reductive group over $\ZZ_p$ such that $\Gscr\otimes_{\ZZ_p}\QQ_p\simeq \mathbf{G}_{\QQ_p}$ and $\Gscr\otimes_{\ZZ_p}\mathbb{F}_p$ is connected. 

We assume that $\gx$ is of abelian type.
For any place $v$ above $p$ in $\mathbf{E}$, Kisin \cite{Kisin-Hodge-Type-Shimura} and Vasiu \cite{Vasiu-Preabelian-integral-canonical-models} constructed a smooth, canonical model $\Sscr_K$  of $\shgx_K$ over $\Ocal_{\mathbf{E}_v}$. Let $\kappa(v)$ denote the residue field of $\Ocal_{\mathbf{E}_v}$, and let $\overline{\FF}_p$ be an algebraic closure of $\kappa(v)$. Set $S_K\colonequals \Sscr_K\otimes_{\Ocal_{\mathbf{E}_v}} \overline{\FF}_p$. We can take a representative $\mu\in \{\mu\}$ defined over $\mathbf{E}_v$, by \cite[Lemma~(1.1.3)(a)]{Kottwitz-Shimura-twisted-orbital}. 
We can also assume that $\mu$ extends to $\mu \colon \GG_{\mathrm{m},\Ocal_{\mathbf{E}_v}}\to \mathscr{G}_{\Ocal_{\mathbf{E}_v}}$ (see \cite[Corollary 3.3.11]{Kim-Rapoport-Zink-uniformization}). It gives rise to a parabolic subgroup $\Pscr\subset \mathscr{G}_{\Ocal_{\mathbf{E}_v}}$, with Levi subgroup $\Lscr$ equal to the centralizer of $\mu$. As explained in \cite[Section~2.5]{Imai-Koskivirta-vector-bundles}, we can assume (after possibly twisting $\mu$) that there is a Borel pair $(\Bscr, \Tscr)$ over $\ZZ_p$ in $\Gscr$ such that $\mathscr{B}_{\Ocal_{\mathbf{E}_v}}\subset \Pscr$. Let $G,P,B,T$ denote the special fibers of $\Gscr, \Pscr, \Bscr, \Tscr$, respectively. By slight abuse of notation, we denote again by $\mu$ its mod $p$ reduction $\mu \colon \GG_{\mathrm{m},k}\to G_k$. 

We define a quotient $\mathbf{G}^c$ of $\mathbf{G}$ by a subtorus of the center of $\mathbf{G}$ as in \cite[Section~2.3]{imai-kato-youcis-prim-real}. We note that $\mathbf{G}^c=\mathbf{G}$ if $\gx$ is of Hodge type by \cite[Remark 2.6]{imai-kato-youcis-prim-real}.  Let $G^c$ be the quotients of $G$ determined by $\mathbf{G}^c$.  Let $\mu^c \colon \GG_{\mathrm{m},k}\to G^c_k$ be the cocharacter induced by $\mu$.  Then $(G^c,\mu^c)$ is a cocharacter datum, and it yields a zip datum as in Section~\ref{subsec-cochar}, where $q=p$.  By \cite[Section~4.1]{Zhang-EO-Hodge} and \cite[Section~3.5]{imai-kato-youcis-prim-real}, there exists a natural smooth morphism
\begin{equation}\label{zeta-Shimura}
\zeta \colon S_K\lra \mathop{\text{$G^c$-{\tt Zip}}}\nolimits^{\mu^c}. 
\end{equation}
This map is also surjective by \cite[Corollary 3.5.3(1)]{Shen-Yu-Zhang-EKOR}. 

Let $T^c$ be the maximal torus of $G^c$ determined by $T$. 
For $\lambda\in X^*_{+,I}(T^c)$, we have a vector bundle $\Vcal_{I}(\lambda)$ on $\mathop{\text{$G^c$-{\tt Zip}}}\nolimits^{\mu^c}$ as in Section~\ref{sec-VB-Gzip}. We denote the pullback of $\Vcal_{I}(\lambda)$ under $\zeta$ by the same symbol. 
The vector bundle $\Vcal_{I}(\lambda)$ on $S_K$ is called \emph{the automorphic vector bundle of weight $\lambda$}.

The flag space of the Siegel modular variety $\Acal_n$ was first introduced by Ekedahl--van der Geer in \cite{Ekedahl-Geer-EO}. It parametrizes pairs $(\underline{A},\Fcal_{\bullet})$, where $\underline{A}=(A,\lambda)\in \Acal_n$ is a principally polarized abelian variety of relative dimension $n$ and $\Fcal_\bullet \subset H^1_{\dR}(A)$ is a full symplectic flag refining the Hodge filtration of $H^1_{\dR}(A)$. In general,  the flag space $\Flag(S_K)$ of $S_K$ was defined in \cite[Section~9.1]{Goldring-Koskivirta-Strata-Hasse} as the fiber product
\begin{equation}
\xymatrix@1@M=7pt{
\Flag(S_K) \ar[d]_{\pi_K} \ar[r]^-{\zeta_{\flag}} & \GF^{\mu} \ar[d]^{\pi} \\
S_K \ar[r]_-{\zeta} & \GZip^\mu\rlap{.}
}
\end{equation}
The stratification $(\Fcal_w)_{w\in W}$ on $\GF^\mu$ induces by pullback via $\zeta_{\flag}$ a stratification $(\Flag(S_K)_w)_{w\in W}$ of $\Flag(S_K)$ by locally closed, smooth subschemes. Moreover, we obtain a line bundle $\Vcal_{\flag}(\lambda)$ on $\Flag(S_K)$. Since $\zeta$ is smooth, pullback and pushforward commute, so we have $\pi_{K,*}(\Vcal_{\flag}(\lambda))=\Vcal_I(\lambda)$. In particular, the space of automorphic forms $H^0(S_K,\Vcal_I(\lambda))$ can be identified with $H^0(\Flag(S_K),\Vcal_{\flag}(\lambda))$.


\setcounter{section}{3}\setcounter{subsection}{4}\setcounter{theorem}{8}\setcounter{equation}{6}

\subsection{Norm of the highest-weight vector}\label{subsec-norm}

Recall that we always have $1\in U_\mu$, where $U_{\mu}\subset G_k$ is the unique open $E$-orbit. Recall the definition of the finite subgroup $L_\varphi\subset L$ given in \eqref{Lphi-equ}. Put $N_\varphi=|L_0(\FF_q)|q^m$, where $L_0$ is the Levi subgroup defined in \eqref{eqL0} and $m\geq 0$ is the smallest integer such that the finite unipotent group $L_\varphi^{\circ}$ is annihilated by $\varphi^m$. For $\lambda,\lambda'\in X^*(T)$ and $f\in V_I(\lambda)$, $f'\in V_I(\lambda')$, let $ff'\in V_I(\lambda+\lambda')$ be the image of $f\otimes f'$ by the map \eqref{Vlambda-natural-map}. For $\lambda\in X^*_{+,I}(T)$ and $f\in V_I(\lambda)$, define
\begin{equation} \label{norm-def-eq}
\Norm_{L_\varphi}(f)\colonequals \left(\prod_{s\in L_0(\FF_q)}s\cdot f \right)^{q^m}\in V_I(N_\varphi \lambda).
\end{equation}
It is clear that $\Norm_{L_\varphi}(f)$ is $L_\varphi$-invariant. In particular, it gives rise to an element in $H^0(\Ucal_\mu,\Vcal(N_\varphi \lambda))$ by Lemma~\ref{lem-Umu-sections}. In general, it is difficult to determine whether $\Norm_{L_\varphi}(f)$ extends to a global section. However, this is possible when $f$ is a highest-weight vector, as we now explain.

Let $f_\lambda\in V_I(\lambda)$ be a nonzero element in the highest-weight line of $V_I(\lambda)$. The following result generalizes \cite[Theorem 2]{Koskivirta-automforms-GZip} (where $P$ was assumed to be defined over $\FF_p$; here we do not make this assumption). For $\alpha\in \Delta^P$, denote by $r_\alpha$ the smallest integer $r\geq 1$ such that $\sigma^r(\alpha)=\alpha$.


\begin{proposition}\label{prop-Norm}
The section $\Norm_{L_\varphi}(f_\lambda)$ extends to a global section over $\GZip^\mu$ if and only if for all $\alpha \in \Delta^P$, the following holds:
\begin{equation}\label{formula-norm}
\sum_{w\in W_{L_0}(\FF_q)} \sum_{i=0}^{r_\alpha-1} q^{i+\ell(w)} \ \left\langle w\lambda, \sigma^i\left(\alpha^\vee\right) \right\rangle\leq 0.
\end{equation}
\end{proposition}

Before giving the proof, we need to recall some facts from \cite[Section~3.1]{Imai-Koskivirta-vector-bundles}. First, we have
\begin{equation}\label{GminusU}
    G_k \setminus U_{\mu} = \bigcup_{\alpha \in \Delta^P} Z_\alpha, \quad Z_\alpha=\overline{E\cdot s_\alpha}, 
\end{equation}
where $E\cdot s_\alpha$ denotes the $E$-orbit of $s_{\alpha}$ and the bar denotes the Zariski closure. This follows easily from Theorem~\ref{thm-E-orb-param}. For any $\alpha\in \Delta^P$, define an open subset 
\begin{equation}
X_\alpha  \colonequals  G_k \setminus \bigcup_{\beta\in \Delta^P,\, \beta\neq \alpha}Z_\beta.
\end{equation}
Then $U_{\mu}\subset X_\alpha$, and $X_{\alpha}\setminus U_{\mu}$ is irreducible. Choose a realization
$(u_{\alpha})_{\alpha \in \Phi}$, and let $\phi_{\alpha}\colon \SL_2\to G$ be the map attached to $\alpha$ (see Section~\ref{subsec-notation}). Set $Y\colonequals E \times \AA^1$ and $Y_0\colonequals E\times \GG_{\mathrm{m}}$. For $\alpha\in \Delta^P$, define $\psi_\alpha \colon Y \to G$ by
\begin{equation}\label{phia}
\psi_\alpha \colon ((x,y),t)\longmapsto x \phi_{\alpha} \left(A(t)\right)y^{-1}, \quad \textrm{where} \ A(t)=\left(
\begin{matrix}
t & 1 \\ -1 & 0 \end{matrix}
\right)\in \SL_{2,k}.
\end{equation}
It satisfies $\psi_\alpha((x,y),t)\in X_\alpha$ for all $((x,y),t)\in Y$ and $\psi_\alpha((x,y),t)\in U_\mu$ if and only if $t\neq 0 $ (see \cite[Proposition 3.1.4]{Imai-Koskivirta-vector-bundles}).

We now prove Proposition~\ref{prop-Norm}. We use a similar argument to that in \cite[Theorem 3.5.3]{Koskivirta-automforms-GZip}. Set $\Ucal'_\mu\colonequals\pi^{-1}(\Ucal_\mu)$, where $\pi\colon \GF^\mu\to \GZip^\mu$ is the natural projection. One  clearly has $\Ucal'_\mu\simeq [E'\backslash U_\mu]$ via the isomorphism $\GF^\mu\simeq [E'\backslash G]$ explained in Section~\ref{subsec-zipflag}. We have an identification
\begin{equation}
    H^0\left(\Ucal_\mu,\Vcal_I\left(N_\varphi\lambda\right)\right)=H^0\left(\Ucal'_\mu,\Vcal_{\flag}\left(N_{\varphi}\lambda\right)\right)
\end{equation}
similarly to \eqref{ident-H0-GF}. In particular, we can view $\Norm_{L_\varphi}(f_{\lambda})$ as a function $h\colon U_\mu\to \AA^1$ satisfying the relation $h(axb^{-1})=\lambda(a)h(x)$ for all $(a,b)\in E'$ and $x\in U_\mu$ (using \eqref{globquot}). Specifically, the function $h$ is given by 
\begin{equation}\label{functionh-eq}
  h\left(x_1x_2^{-1}\right)=\Norm_{L_\varphi}(f_\lambda)\left(\theta_L^P(x_1)^{-1}\right)
\end{equation}
for all $(x_1,x_2)\in E$ using \eqref{eq:H0expid}. 
The function $h$ is well defined because $\Norm_{L_\varphi}(f_\lambda)$ is $L_\varphi$-invariant. Furthermore, $\Norm_{L_\varphi}(f_\lambda)$ extends to $\GZip^\mu$ if and only if $h$ extends to a function $G\to \AA^1$. By the strategy explained in \cite[Section~3.2]{Koskivirta-automforms-GZip} and in \cite[Section~3.1]{Imai-Koskivirta-vector-bundles}, the function $h$ extends to $G$ if and only if for each $\alpha\in \Delta^P$, the function $h\circ \psi_\alpha \colon Y_0\to \AA^1$ extends to a function $Y\to \AA^1$. It remains to compute the $t$-valuation of the function $h\circ \psi_\alpha$, viewed as an element of $R[t,\frac{1}{t}]$, where $R=k[E]$ is the ring of functions of $E$. Put 
\begin{equation}\label{malpha-equ}
 m_{\alpha} =\min \left\{ m \geq 1 \mid 
 \sigma^{-m}(\alpha) \notin I \right\}, \quad \alpha \in \Delta^P
\end{equation}
and $t_{\alpha}=t^{-1}\alpha(\varphi(\delta_{\alpha}(t)))^{-1} =t \alpha(\delta_{\alpha}(t))^{-1} \in t^{\QQ}$, where $t$ is an indeterminate and $\delta_{\alpha}=\wp_*^{-1}(\alpha^{\vee})$ as defined in Section~\ref{subsec-global-sections}. Set
\[
 u_{t,\alpha}=\prod_{i=1}^{m_{\alpha}-1} \phi_{\sigma^{-i}(\alpha)} 
 \left( \begin{pmatrix}
1&-t_{\alpha}^{\frac{1}{q^i}}\\0&1
\end{pmatrix} \right),  
\]
where the product is taken in increasing order of indices. By the proof of \cite[Proposition 3.1.4]{Imai-Koskivirta-vector-bundles}, for all $(x,y)\in E$ and $t\in \GG_{\mathrm{m}}$, we can write $\psi_\alpha((x,y),t)=x_1x_2^{-1}$ with $(x_1,x_2)\in E$ and
\begin{equation}
    x_1=x\phi_{\alpha} \left(\begin{pmatrix}
1&0\\-t^{-1}&1
\end{pmatrix} \right) \delta_{\alpha}(t) 
u_{t,\alpha}.
\end{equation}
By the definition of $m_\alpha$, all the roots $\sigma^{-i}(\alpha)$ (for $1\leq i\leq m_\alpha-1$) appearing in the formula of $u_{t,\alpha}$ lie in $I$. Using \eqref{functionh-eq}, we deduce 
\begin{align}
h\circ \psi_\alpha((x,y),t)&=\Norm_{L_\varphi}(f_\lambda)\left(u_{t,\alpha}^{-1}\delta_\alpha(t)^{-1}\theta_L^P(x)^{-1}\right) \\
&=\left(\prod_{s\in L_0(\FF_q)} f_\lambda\left(su_{t,\alpha}^{-1}\delta_\alpha(t)^{-1}\theta_L^P(x)^{-1}\right) \right)^{q^m}.
\end{align}
Consider the element $f_\lambda(su_{t,\alpha}^{-1}\delta_\alpha(t)^{-1}\theta_L^P(x)^{-1})$, which lies in $R[t^{\QQ}]$. We can still speak of the $t$-valuation of this element, which is a rational number. Equivalently, to simplify notation, we change $\theta_L^P(x)^{-1}$ to a generic element $g\in L$, and we compute the $t$-valuation of $F_s(t,g)\colonequals f_\lambda(su_{t,\alpha}^{-1}\delta_\alpha(t)^{-1}g)$, viewed as an element of $k[L][t^\QQ]$. Let $v_\alpha(s)$ be this valuation. 
We put $B^+_L=B^+ \cap L$. 
Define a parabolic subgroup of $L$ by $Q_0\colonequals L_0 B^+_L$. It is clear that $u_{t,\alpha}$ lies in $R_{\mathrm{u}}(Q_0)$; thus for all $s\in L_0(\FF_q)$, we have $su_{t,\alpha}^{-1}s^{-1}\in R_{\mathrm{u}}(Q_0)$. Since $f_\lambda$ is invariant by $R_{\mathrm{u}}(B^+_L)$, we obtain 
$F_s(t,g)=f_\lambda(s\delta_\alpha(t)^{-1}g)$. Now, the rest of the proof is completely similar to that of \cite[Theorem 3.5.3]{Koskivirta-automforms-GZip}. We recall it briefly.

Let $B_{L_0}\colonequals B\cap L_0$ and $B^+_{L_0}\colonequals B^+\cap L_0$. If we change $s$ to $bs$ with $b\in B^{+}_{L_0}(\FF_q)$, then $v_\alpha(bs)=v_\alpha(s)$. Indeed, this follows from $f_\lambda(bs\delta_\alpha(t)^{-1}g)=\lambda(b)^{-1} f_\lambda(s\delta_\alpha(t)^{-1}g)$ since $f_\lambda$ is a $B^+_L$-eigenfunction. Similarly, we claim that $v_\alpha(sb)=v_\alpha(s)$ for all $b\in B^+_{L_0}(\FF_q)$. By symmetry, it suffices to show $v_\alpha(sb)\geq v_\alpha(s)$. We can write
$F_{sb}(t,g)=F_s(t,\Gamma(t)g)$, where
\begin{equation}
    \Gamma(t)\colonequals
    \delta_\alpha(t) b \delta_\alpha(t)^{-1}.
\end{equation}
We view $\Gamma$ as a map $\Gamma\colon \spec(k[t^\QQ])\to L$. Since $\alpha\in \Delta^P$, the cocharacter $\alpha^\vee$ is anti-$L$-dominant. It follows that for all $j\in \ZZ$, $\sigma^j(\alpha^\vee)$ is an anti-$L_0$-dominant quasi-cocharacter. It is easy to see that $\delta_\alpha$ is explicitly given by the formula
\begin{equation}\label{eq-deltaalpha}
    \delta_\alpha=-\frac{1}{q^{r_\alpha}-1}\sum_{j=0}^{r_\alpha-1}q^j\sigma^j\left(\alpha^\vee\right). 
\end{equation}
In particular, $\delta_\alpha$ is $L_0$-dominant. We deduce that the function $\Gamma$ extends to a map $\spec(k[t^{\QQ_{\geq 0}}])\to L$. This follows simply from the fact that $\delta_\alpha(t)$ acts on the root space $U_\beta$ (for $\beta\in \Phi$) by $t^{\langle \beta,\delta_\alpha \rangle}$, using \eqref{eq:phiconj}. Write $F_s(t,g)=t^{v_\alpha(s)}F_{s,0}(t,g)$, where $F_{s,0}(t,g)$ is an element of $k[L][t^\QQ]$ whose $t$-valuation is $0$. Then $F_{sb}(t,g)=t^{v_\alpha(s)}F_{s,0}(t,\Gamma(t)g)$ and 
$F_{s,0}(t,\Gamma(t)g) \in k[L][t^{\QQ_{\geq 0}}]$. 
Hence $v_\alpha(sb)\geq v_\alpha(s)$ as claimed.

Now, consider the Bruhat decomposition of $L_0(\FF_q)$:
\begin{equation}
    L_0\left(\FF_q\right)=\bigsqcup_{w\in W_{L_0}(\FF_q)} B^+_{L_0}\left(\FF_q\right) w B^+_{L_0}\left(\FF_q\right)
\end{equation}
as in \cite[Lemma 3.4.4]{Koskivirta-automforms-GZip}. By \cite[Lemma 3.4.5]{Koskivirta-automforms-GZip}, one has
\begin{equation}
  \left|B^+_{L_0}\left(\FF_q\right) w B^+_{L_0}\left(\FF_q\right)\right| =  \left|T\left(\FF_q\right)\right|q^{\dim(R_{\mathrm{u}}(B_{L_0}))+\ell(w)}.
\end{equation} 
Thus, we can determine $v_\alpha$ completely from the values $v_\alpha(w)$ for $w\in W_{L_0}(\FF_q)$. Similarly to \cite[Proposition 3.5.2]{Koskivirta-automforms-GZip}, we have $v_\alpha(w)=\langle w\lambda, \delta_\alpha \rangle$. We deduce that the $t$-valuation of $h\circ \psi_\alpha((x,y),t)$ is
\begin{equation}
 q^m\sum_{s\in L_0(\FF_q)} v_\alpha(s) =   q^m \left|T\left(\FF_q\right)\right|q^{\dim(R_{\mathrm{u}}(B_{L_0}))} \sum_{w\in W_{L_0}(\FF_q)} q^{\ell(w)}\langle w\lambda, \delta_\alpha \rangle.
\end{equation}
The statement of Proposition~\ref{prop-Norm} then follows by replacing $\delta_\alpha$ by the expression in \eqref{eq-deltaalpha}.

\phantomsection
\newcommand{\etalchar}[1]{$^{#1}$}
\providecommand{\bysame}{\leavevmode\hbox to3em{\hrulefill}\thinspace}

\begin{thebibliography}{EvdG09+++}

\bibitem[And23]{Andreatta-modp-period-maps}
F.~Andreatta, \emph{On two {${\rm mod}\,p$} period maps: {E}kedahl--{O}ort and
  fine {D}eligne--{L}usztig stratifications}, Math.\ Ann.\ \textbf{385} (2023),
no.~1-2, 511--550, \doi{10.1007/s00208-021-02356-7}.

\bibitem[AG05]{Andreatta-Goren-book}
F.~Andreatta and E.\,Z.~Goren, \emph{Hilbert Modular Forms: mod $p$ and $p$-Adic
  Aspects}, Mem.\ Amer.\ Math.\ Soc.\ \textbf{173} (2005), no.~819, 
\doi{10.1090/memo/0819}.

  
\bibitem[ABD{\etalchar{+}}11]{SGA3}
M.~Artin, J.-E.~Bertin, M.~Demazure, A.~Grothendieck, P.~Gabriel, M.~Raynaud, and J.-P.~Serre,  \emph{Sch\'emas en groupes (SGA 3). Tome III. Structure des sch\'emas en groupes r\'eductifs}, S\'eminaire de G\'eom\'etrie Alg\'ebrique du Bois Marie 1962--64, Revised and annotated edition of the 1970 French original, edited by P.~Gille and P.~Polo, Doc.\ Math. (Paris), vol.~7, Soc.\ Math.\ France,
Paris, 2011.
%M.~Artin, J.-E.~Bertin, M.~Demazure, P.~Gabriel, A.~Grothendieck, M.~Raynaud and J.-P.~Serre, \emph{Sch\'emas en groupes. Fasc.~6: Expos\'es 19 \`a 22  et Fasc.~7: Expos\'es 23 \`a 26},S\'eminaire de G\'eom\'etrie Alg\'ebrique de l'Inst.\ des Hautes \'Etudes Sci., 1963/64, Inst.\ Hautes \'Etudes Sci., Paris, 1965.

\bibitem[Bou68]{bourbaki-lie-4-6}
  N.~Bourbaki, \emph{\'El\'ements de math\'ematique. Fasc. XXXIV. Groupes et alg\`ebres de Lie. Chapitres IV--VI}, Actualit\'es Sci.\ Indust., No.~1337, Hermann, Paris, 1968. 


\bibitem[Del79]{Deligne-Shimura-varieties}
P.~Deligne, \emph{Vari\'et\'es de {S}himura: {I}nterpr\'etation modulaire, et
  techniques de construction de mod\`eles canoniques}, in: \emph{Automorphic forms,  representations and {$L$}-functions. Part~2} ({O}regon {S}tate {U}niv., {C}orvallis, Ore., 1977), pp.~247--289, {P}roc.\ {S}ympos.\ {P}ure {M}ath., vol.~33, Amer.\ Math.\ Soc., Providence, RI,
  1979, \doi{10.1090/pspum/033.2/546620}.

\bibitem[DK17]{Diamond-Kassaei-comp-minimal}
F.~Diamond and P.~Kassaei, \emph{Minimal weights of {H}ilbert modular forms in
  characteristic {$p$}}, Compos.\ Math.\ \textbf{153} (2017), no.~9, 1769--1778,
\doi{10.1112/S0010437X17007230}.

\bibitem[DK23]{Diamond-Kassaei-cone-minimal}
\bysame, \emph{The {C}one of {M}inimal {W}eights for {M}od $p$ {H}ilbert
  {M}odular {F}orms}, Int.\ Math.\ Res.\ Not.\ IMRN \textbf{2023}, no.~14, 12148--12171, 
\doi{10.1093/imrn/rnac121}.

\bibitem[EvdG09]{Ekedahl-Geer-EO}
T.~Ekedahl and G.~van~der Geer, \emph{Cycle Classes of the {E}-{O}
  Stratification on the Moduli of Abelian Varieties}, in: \emph{Algebra, Arithmetic and  Geometry: In Honor of Yu.\,I.~Manin. Vol.~1}, pp.~567--636, 
  Progr.\ Math., vol.~269, Birkh\"auser Boston, Boston, MA, 2009, 
\doi{10.1007/978-0-8176-4745-2_13}.

\bibitem[GK18]{Goldring-Koskivirta-global-sections-compositio}
W.~Goldring and J.\,S.~Koskivirta, \emph{Automorphic vector bundles with global
  sections on {$G$}-{\tt Zip}$^{{\mathcal{Z}}}$-schemes}, Compos.\ Math.\
  \textbf{154} (2018), no.~12, 2586--2605,
  \doi{10.1112/S0010437X18007467}.
  
\bibitem[GK19a]{Goldring-Koskivirta-Strata-Hasse}
\bysame, \emph{Strata {H}asse invariants, {H}ecke algebras and {G}alois
  representations}, Invent.\ Math.\ \textbf{217} (2019), no.~3, 887--984, 
\doi{10.1007/s00222-019-00882-5}.

\bibitem[GK19b]{Goldring-Koskivirta-zip-flags}
\bysame, \emph{Stratifications of Flag Spaces and Functoriality}, Int.\ Math.\ Res.\ Not.\ IMRN
\textbf{2019}, no.~12, 3646--3682, \doi{10.1093/imrn/rnx229}.

\bibitem[GK22a]{Goldring-Koskivirta-divisibility}
\bysame, \emph{Divisibility of mod $p$ automorphic forms and the cone
conjecture for certain {S}himura varieties of {H}odge-type}, preprint \arXiv{2211.16817} (2022). 

\bibitem[GK22b]{Goldring-Koskivirta-GS-cone}
\bysame, \emph{Griffiths-{S}chmid conditions for automorphic forms via
characteristic $p$}, preprint
\href{https://arxiv.org/abs/2211.16819}{{arXiv:\allowbreak 2211.16819}} (2022).


\bibitem[GS69]{Griffiths-Schmid-homogeneous-complex-manifolds}
P.~Griffiths and W.~Schmid, \emph{Locally homogeneous complex manifolds}, Acta.\
  Math.\ \textbf{123} (1969), 253--302, 
\doi{10.1007/BF02392390}.

\bibitem[IKY23]{imai-kato-youcis-prim-real}
N.~Imai, H.~Kato and A.~Youcis, \emph{The prismatic realization functor for
shimura varieties of abelian type}, preprint \arXiv{2310.08472} (2023).

\bibitem[IK21]{Imai-Koskivirta-vector-bundles}
N.~Imai and J.\,S.~Koskivirta, \emph{Automorphic vector bundles on the stack of
  $G$-zips}, Forum Math.\ Sigma \textbf{9} (2021), Paper No.~e37, 
\doi{10.1017/fms.2021.32}.

\bibitem[IK24]{Imai-Koskivirta-partial-Hasse}
\bysame, \emph{Partial {H}asse invariants for {S}himura varieties of
  {H}odge-type}, Adv.\ Math.\ \textbf{440} (2024), Paper No.~109518,
\doi{10.1016/j.aim.2024.109518}.

\bibitem[Jan03]{jantzen-representations}
J.\,C.~Jantzen, \emph{Representations of algebraic groups}, 2nd ed., Math.\ Surveys  Monogr., vol.~107, Amer.\ Math.\ Soc., Providence, RI,
  2003.

\bibitem[Kim18]{Kim-Rapoport-Zink-uniformization}
W.~Kim, \emph{Rapoport--{Z}ink uniformization of {S}himura varieties}, Forum
  Math.\ Sigma \textbf{6} (2018), Paper No.~e16,  
\doi{10.1017/fms.2018.18}.

\bibitem[Kis10]{Kisin-Hodge-Type-Shimura}
M.~Kisin, \emph{Integral models for {S}himura varieties of abelian type},
  J.~Amer.\ Math.\ Soc.\ \textbf{23} (2010), no.~4, 967--1012, 
\doi{10.1090/S0894-0347-10-00667-3}.

\bibitem[Kos19]{Koskivirta-automforms-GZip}
J.-S.~Koskivirta, \emph{Automorphic Forms on the Stack of {$G$}-Zips}, Results
  Math.\ \textbf{74} (2019), no.~3, Paper No.~91,  
\doi{10.1007/s00025-019-1021-z}.

\bibitem[KW18]{Koskivirta-Wedhorn-Hasse}
J.-S.~Koskivirta and T.~Wedhorn, \emph{Generalized {$\mu$}-ordinary {H}asse
  invariants}, J.~Algebra \textbf{502} (2018), 98--119, 
\doi{10.1016/j.jalgebra.2018.01.011}.

\bibitem[Kot86]{Kottwitz-elliptic-singular}
R.\,E.~Kottwitz, \emph{Stable trace formula: Elliptic singular terms}, Math.\ Ann.\
  \textbf{275} (1986), no.~3, 365--399, \doi{10.1007/BF01458611}.

\bibitem[Kot84]{Kottwitz-Shimura-twisted-orbital}
\bysame, \emph{Shimura varieties and twisted orbital integrals},
  Math.\ Ann.\ \textbf{269} (1984), no.~3, 287--300. 
  \doi{10.1007/BF01450697}.
  

\bibitem[KS24]{Kret-Shin-GSO}
A.~Kret and S.\,W.~Shin, \emph{Galois representations for even general special
  orthogonal groups}, J.~Inst.\ Math.\ Jussieu \textbf{23} (2024), no.~5,
  1959--2050, \doi{10.1017/S1474748023000427}.

\bibitem[LS18]{Lan-Stroh-stratifications-compactifications}
K.-W.~Lan and B.~Stroh, \emph{Compactifications of subschemes of integral models  of {S}himura varieties}, Forum Math.\ Sigma \textbf{6} (2018), Paper No.~e18, 
\doi{10.1017/fms.2018.20}.

\bibitem[MW04]{Moonen-Wedhorn-Discrete-Invariants}
B.~Moonen and T.~Wedhorn, \emph{Discrete invariants of varieties in positive
  characteristic}, Int.\ Math.\ Res.\ Not.\ IMRN \textbf{2004}, no.~72, 3855--3903, \doi{10.1155/S1073792804141263}.

\bibitem[PWZ11]{Pink-Wedhorn-Ziegler-zip-data}
R.~Pink, T.~Wedhorn and P.~Ziegler, \emph{Algebraic zip data}, Doc.\ Math.\
  \textbf{16} (2011), 253--300, 
\doi{10.4171/dm/332}.

\bibitem[PWZ15]{Pink-Wedhorn-Ziegler-F-Zips-additional-structure}
\bysame, \emph{${F}$-zips with additional structure}, Pacific J.~Math.\
  \textbf{274} (2015), no.~1, 183--236, 
\doi{10.2140/pjm.2015.274.183}.

\bibitem[RR85]{Ramanan-Ramanathan-projective-normality}
S.~Ramanan and A.~Ramanathan, \emph{Projective normality of flag varieties and
  {S}chubert varieties}, Invent.\ Math.\ \textbf{79} (1985), no.~2, 217--224,
\doi{10.1007/BF01388970}.

\bibitem[SYZ21]{Shen-Yu-Zhang-EKOR}
X.~Shen, C.-F.~Yu and C.~Zhang, \emph{E{KOR} strata for {S}himura varieties
  with parahoric level structure}, Duke Math.~J.\ \textbf{170} (2021), no.~14,
  3111--3236, 
\doi{10.1215/00127094-2021-0047}.

\bibitem[Spr98]{Springer-Linear-Algebraic-Groups-book}
T.\,A.~Springer, \emph{Linear Algebraic Groups}, 2nd ed., Progr.\ Math.,
  vol.~9, Birkh\"auser Boston, Boston, MA, 1998, 
\doi{10.1007/978-0-8176-4840-4}.

\bibitem[Vas99]{Vasiu-Preabelian-integral-canonical-models}
A.~Vasiu, \emph{Integral canonical models of {S}himura varieties of preabelian
  type}, Asian J.~Math.\ \textbf{3} (1999), no.~2, 401--518,
\doi{10.4310/AJM.1999.v3.n2.a8}.

\bibitem[Wed14]{Wedhorn-bruhat}
T.~Wedhorn, \emph{Bruhat strata and ${F}$-zips with additional structure},
  M\"unster J.~Math.\ \textbf{7} (2014), no.~2, 529--556, \doi{10.17879/58269760517}. 

\bibitem[WZ23]{Wedhorn-Ziegler-tautological}
T.~Wedhorn and P.~Ziegler, \emph{Tautological rings of {S}himura varieties and
  cycle classes of {E}kedahl--{O}ort strata}, Algebra Number Theory \textbf{17}
  (2023), no.~4, 923--980, \doi{10.2140/ant.2023.17.923}.

\bibitem[XZ19]{Xiao-Zhu-on-vector-valued}
L.~Xiao and X.~Zhu, \emph{On vector-valued twisted conjugation invariant
  functions on a group} (with an appendix by S.~Donkin), in: \emph{Representations of Reductive Groups}, pp.~361--425, Proc.\ Sympos.\
  Pure Math., vol.~101, Amer.\ Math.\ Soc., Providence, RI, 2019,  
\doi{10.1090/pspum/101/14}.

\bibitem[Zha18]{Zhang-EO-Hodge}
C.~Zhang, \emph{Ekedahl-{O}ort Strata for Good Reductions of {S}himura
  Varieties of {H}odge Type}, Canad.~J.\ Math.\ \textbf{70} (2018), no.~2,
  451--480,
\doi{10.4153/CJM-2017-020-5}.

\end{thebibliography}
\end{document}
